\section*{अध्याय \ref{ch:g10:concentration-and-dilution} --- सांद्रता और तनुकरण}

\begin{solution}{exo:g10:concentration-and-dilution:1}
$C_m = \qty{5.0}{g} / \qty{0.250}{L} = \qty{20}{g/L}$।
\end{solution}

\begin{solution}{exo:g10:concentration-and-dilution:2}
$c = \qty{0.050}{mol} / \qty{0.200}{L} = \qty{0.25}{mol/L}$।
\end{solution}

\begin{solution}{exo:g10:concentration-and-dilution:3}
$M(\ce{CuSO4}) = 63.5 + 32.1 + 4 \times 16.0 = \qty{159.6}{g/mol}$;
$n = 0.10 \times 0.1000 = \qty{0.0100}{mol}$; $m = 0.0100 \times 159.6 =
\qty{1.60}{g}$।
\end{solution}

\begin{solution}{exo:g10:concentration-and-dilution:4}
$F = 1.0 / 0.050 = 20$; $V_0 = \qty{100.0}{mL} / 20 = \qty{5.0}{mL}$।
\end{solution}

\begin{solution}{exo:g10:concentration-and-dilution:5}
$C_m = c \times M = 0.10 \times 58.5 = \qty{5.85}{g/L}$, लगभग
\qty{5.9}{g/L}।
\end{solution}

\begin{solution}{exo:g10:concentration-and-dilution:6}
$V_0 = c_1 V_1 / c_0 = 0.020 \times 250.0 / 0.50 = \qty{10.0}{mL}$
(\omterm{def:g10:concentration-and-dilution:dilution}{तनुकरण गुणांक} 25)। एक \qty{10.0}{mL} का आयतनमितीय पिपेट और एक
\qty{250.0}{mL} का आयतनमितीय फ़्लास्क।
\end{solution}

\begin{solution}{exo:g10:concentration-and-dilution:7}
आयतनमितीय फ़्लास्क पर एक ही निशान होता है, जो एक आयतन के लिए बना है और एक
प्रतिशत के छोटे-से अंश तक यथार्थ है; बीकर के अंशांकन केवल मोटा अंदाज़ा देते हैं।
इसी तरह आयतनमितीय पिपेट एक आयतन बहुत यथार्थता से देता है, जबकि मापक सिलिंडर
चौड़ा होता है और कम परिशुद्धता से
पढ़ा जाता है।
\end{solution}

\begin{solution}{exo:g10:concentration-and-dilution:8}
\qty{0.02}{mol/L} और \qty{0.04}{mol/L} के बीच: लगभग \qty{0.03}{mol/L}।
अधिक परिशुद्धता के लिए इन दोनों के बीच \omterm{def:g10:concentration-and-dilution:standard}{मानक विलयन} जोड़िए, जैसे
0.025, 0.030 और \qty{0.035}{mol/L} पर: जहाँ ज़रूरत है, वहाँ बारीक श्रेणी।
\end{solution}

\begin{solution}{exo:g10:concentration-and-dilution:9}
$M(\ce{C6H12O6}) = \qty{180.0}{g/mol}$; $c = 50 / 180.0 =
\qty{0.28}{mol/L}$।
\end{solution}

\begin{solution}{exo:g10:concentration-and-dilution:10}
$M = \qty{342.0}{g/mol}$, इसलिए $n = \qty{1.00}{mol}$ और $c =
\qty{1.00}{mol/L}$। दस गुना तनु करने पर: $c = \qty{0.100}{mol/L}$, यानी
$C_m = 0.100 \times 342.0 = \qty{34.2}{g/L}$।
\qty{0.020}{L} की एक चम्मच में $34.2 \times 0.020 = \qty{0.68}{g}$ सुक्रोस है।
\end{solution}

\begin{solution}{exo:g10:concentration-and-dilution:11}
(क) \omterm{def:g6:solutions-and-solubility:solute}{विलेय} का वही द्रव्यमान दुगने आयतन में: \qty{6.0}{g/L}।
(ख) आयतन लगभग आधा हो जाता है, \omterm{def:g6:solutions-and-solubility:solute}{विलेय} का द्रव्यमान वही रहता है: लगभग
\qty{24}{g/L}।
\end{solution}

\begin{solution}{exo:g10:concentration-and-dilution:12}
\begin{enumerate}
  \item 10 के तीन \omterm{def:g10:concentration-and-dilution:dilution}{तनुकरण} मिलकर $10^3$ का \omterm{def:g10:concentration-and-dilution:dilution}{तनुकरण} बनाते हैं:
        $\qty{0.10}{mol/L} / 1000 = \qty{1.0e-4}{mol/L}$।
  \item 1000 गुना के एक ही \omterm{def:g10:concentration-and-dilution:dilution}{तनुकरण} के लिए या तो \omterm{def:g10:concentration-and-dilution:dilution}{मूल विलयन} का बहुत छोटा आयतन
        (\qty{100}{mL} में \qty{0.1}{mL}) चाहिए, जिसे पिपेट से यथार्थता से
        मापना असंभव है, या बहुत बड़ा फ़्लास्क (\qty{1}{L} में \qty{1}{mL})।
        तीन दस-गुना तनुकरणों में साधारण पिपेट और
        फ़्लास्क काम आते हैं।
\end{enumerate}
\end{solution}

\begin{solution}{exo:g10:concentration-and-dilution:13}
\begin{enumerate}
  \item $M(\ce{CaCl2}) = 40.1 + 2 \times 35.5 = \qty{111.1}{g/mol}$;
        $n = 11.1 / 111.1 = \qty{0.100}{mol}$; $c = 0.100 / 0.5000 =
        \qty{0.200}{mol/L}$।
  \item हर \ce{CaCl2} एक \ce{Ca^2+} और दो \ce{Cl-} देता है:
        $[\ce{Ca^2+}] = \qty{0.200}{mol/L}$ और $[\ce{Cl-}] =
        \qty{0.400}{mol/L}$।
\end{enumerate}
\end{solution}

\begin{solution}{exo:g10:concentration-and-dilution:14}
\begin{enumerate}
  \item $r = \qty{0.60}{cm}$, $h = \qty{0.20}{cm}$:
        $V = \pi \times 0.60^2 \times 0.20 = \qty{0.23}{cm^3}$, यानी
        \qty{0.23}{mL} पानी ज़्यादा।
  \item वही \omterm{def:g6:solutions-and-solubility:solute}{विलेय} \qty{100.00}{mL} के बजाय \qty{100.23}{mL} में है:
        सांद्रता $0.23 / 100.23 \approx
        \qty{0.23}{\percent}$ कम है।
\end{enumerate}
\end{solution}

\begin{solution}{exo:g10:concentration-and-dilution:15}
$n = 0.20 \times 0.100 + 0.10 \times 0.300 = 0.020 + 0.030 =
\qty{0.050}{mol}$, \qty{0.400}{L} में: $c = \qty{0.125}{mol/L}$। यह सीधा औसत,
\qty{0.15}{mol/L}, नहीं है: तनु विलयन तीन गुना अधिक है, इसलिए \omterm{def:g6:pure-substances-and-mixtures:mixture}{मिश्रण}
\qty{0.10}{mol/L} के अधिक क़रीब है। यह आयतनों से भारित
औसत है।
\end{solution}

\begin{solution}{pb:g10:concentration-and-dilution:1}
\textbf{1.} \qty{0.100}{L} में \qty{0.9}{g}: $C_m = \qty{9.0}{g/L}$।

\textbf{2.} \qty{1000}{mL} में \qty{9.0}{g}, यानी प्रति मिलीलीटर
\qty{9.0}{mg}।

\textbf{3.} $m = 9.0 \times 0.500 = \qty{4.5}{g}$।

\textbf{4.} हाँ: प्रति लीटर लगभग \qty{9}{g}, जबकि संतृप्ति पर प्रति किलोग्राम पानी
लगभग \qty{360}{g}, यानी लगभग चालीस गुना कम।

\textbf{5.} $9.0 \times 1.0 = \qty{9.0}{g}$ नमक।

\textbf{6.} $M(\ce{NaCl}) = 23.0 + 35.5 = \qty{58.5}{g/mol}$।

\textbf{7.} $c = 9.0 / 58.5 = \qty{0.154}{mol/L}$।

\textbf{8.} $n = 0.154 \times 0.500 = \qty{0.0769}{mol}$ (या
$4.5 / 58.5$)।

\textbf{9.} हर \ce{NaCl} एक \ce{Na+} और एक \ce{Cl-} देता है: दोनों
\qty{0.154}{mol/L} पर।

\textbf{10.} $0.0769 \times \num{6.02e23} = \num{4.63e22}$ सोडियम \omterm{def:g9:ions:ion}{आयन}।

\textbf{11.} $F = 200 / 9.0 = 22.2$।

\textbf{12.} थैली का \qty{4.5}{g} पूरा का पूरा सांद्र विलयन से आता है:
$V_0 = \qty{4.5}{g} / \qty{200}{g/L} = \qty{0.0225}{L} =
\qty{22.5}{mL}$।

\textbf{13.} $V_0 = V_1 / F = 500 / 22.2 = \qty{22.5}{mL}$।

\textbf{14.} $c_0 = 200 / 58.5 = \qty{3.42}{mol/L}$;
$c_0 V_0 = 3.42 \times 0.0225 = \qty{0.0769}{mol}$ और
$c_1 V_1 = 0.154 \times 0.500 = \qty{0.0769}{mol}$: बराबर।

\textbf{15.} लगभग $500 - 22.5 = \qty{478}{mL}$ जीवाणुरहित पानी।

\textbf{16.} \qty{25}{mL} का एक अंशांकित पिपेट (या एक ब्यूरेट), जिसे
\qty{0.1}{mL} तक पढ़ा जाए।

\textbf{17.} \qty{500.0}{mL} के आयतनमितीय फ़्लास्क में, निशान तक पूरा करके।

\textbf{18.} $C_m = 200 \times 23.0 / 500 = \qty{9.20}{g/L}$, यानी
$0.20 / 9.0 \approx \qty{2.2}{\percent}$ अधिक सांद्र।

\textbf{19.} \qty{500}{mL} की एक थैली के लिए सांद्र (``20 \%'') विलयन का
\qty{22.5}{mL}।
\end{solution}
